\documentclass[11pt,a4paper,fontset=fandol]{ctexart}

\usepackage[a4paper,top=2.35cm,bottom=2.55cm,left=2.3cm,right=2.3cm,headheight=18pt,headsep=0.55cm,footskip=26pt]{geometry}
\usepackage{amsmath,amssymb,amsthm}
\usepackage{fancyhdr}
\usepackage{lastpage}
\usepackage{enumitem}
\usepackage{titlesec}
\usepackage{xcolor}
\usepackage{tikz}
\usepackage{hyperref}
\usepackage{bookmark}
\usepackage[most]{tcolorbox}
\usepackage{setspace}
\usepackage{array}
\usetikzlibrary{arrows.meta,positioning}

\hypersetup{
  colorlinks=true,
  linkcolor=blue!50!black,
  urlcolor=blue!50!black,
  citecolor=blue!50!black
}

\setstretch{1.15}
\setlength{\parindent}{2em}
\setlength{\parskip}{0.28em}
\allowdisplaybreaks
\setlist[enumerate]{itemsep=0.35em, topsep=0.35em, leftmargin=2.2em}
\setlist[itemize]{itemsep=0.25em, topsep=0.25em, leftmargin=1.9em}

\definecolor{mainblue}{RGB}{36,83,140}
\definecolor{lightblue}{RGB}{241,246,252}
\definecolor{lightgray}{RGB}{246,246,246}
\definecolor{accent}{RGB}{153,76,0}
\definecolor{rulegray}{RGB}{110,118,128}

\titleformat{\section}
  {\Large\bfseries\color{mainblue}}
  {\thesection.}
  {0.45em}
  {}
  [\vspace{0.25em}\color{rulegray}\titlerule]
\titlespacing*{\section}{0pt}{1.2em}{0.8em}
\titleformat{\subsection}{\large\bfseries\color{mainblue}}{\thesubsection}{0.5em}{}

\pagestyle{fancy}
\fancyhf{}
\fancyhead[L]{\small 函数图像对称提高训练题单}
\fancyhead[C]{\small 代数专题：图像对称与方程变形}
\fancyhead[R]{\small 刘文博}
\fancyfoot[C]{\small 第 \thepage 页 / 共 \pageref{LastPage} 页}
\renewcommand{\headrulewidth}{0.55pt}
\renewcommand{\footrulewidth}{0.35pt}

\newtcolorbox{goalbox}[1]{
  breakable,
  colback=lightblue,
  colframe=mainblue,
  boxrule=0.7pt,
  arc=1mm,
  title=\textbf{#1},
  fonttitle=\bfseries
}

\newtcolorbox{notebox}[1]{
  breakable,
  colback=lightgray,
  colframe=black!50,
  boxrule=0.55pt,
  arc=1mm,
  title=\textbf{#1},
  fonttitle=\bfseries
}

\newtheoremstyle{formalexample}
  {0.8em}{0.8em}{\normalfont}{}{\bfseries\color{accent}}{．}{0.6em}{}
\theoremstyle{formalexample}
\newtheorem{exercise}{题目}[section]
\newenvironment{solution}{\par\noindent\textbf{解：}}{\hfill$\square$\par}

\begin{document}

\thispagestyle{empty}
\begin{titlepage}
\centering
\vspace*{0.9cm}
\begin{tcolorbox}[
  enhanced,
  width=0.92\textwidth,
  colback=white,
  colframe=mainblue,
  boxrule=1pt,
  arc=0mm,
  left=6mm,
  right=6mm,
  top=8mm,
  bottom=8mm
]
\centering
{\fontsize{24}{30}\selectfont\bfseries 函数图像对称提高训练题单\par}
\vspace{0.6cm}
{\fontsize{17}{22}\selectfont 适合初中阶段函数图像专题巩固训练\par}

\vspace{1cm}
\begin{tcolorbox}[colback=lightblue,colframe=mainblue,width=0.84\textwidth,boxrule=1pt]
\centering
{\large 训练目标：把“会套公式”提升到“会判断、会解释、会综合变换”}\\[0.35em]
{\normalsize 建议分 2--3 次完成，先独立做题，再对照详细解答订正}
\end{tcolorbox}

\vspace{1cm}
\renewcommand{\arraystretch}{1.42}
\begin{tabular}{>{\bfseries}p{3.5cm}p{8cm}}
题目数量： & 16 题，按层次递进 \\
专题重点： & 关于 x 轴、关于 y 轴、关于原点对称；奇偶性判断；对称与平移衔接 \\
答案形式： & 末尾附详细解答，强调“为什么这样改方程” \\
编译方式： & XeLaTeX \\
\end{tabular}

\vspace{0.9cm}
{\large \today\par}
\end{tcolorbox}
\end{titlepage}

\tableofcontents
\newpage

\section{使用建议}

\begin{goalbox}{做函数对称题时，优先问自己的 4 个问题}
\begin{itemize}
  \item 这是关于 x 轴、关于 y 轴，还是关于原点对称？
  \item 横坐标变了吗？纵坐标变了吗？
  \item 这是在函数外面加负号，还是在函数里面把 $x$ 换成 $-x$？
  \item 我能不能用特殊点、顶点或偶函数/奇函数性质来检验结果？
\end{itemize}
\end{goalbox}

\begin{notebox}{建议计时}
\begin{itemize}
  \item 基础题：每题 2--4 分钟；
  \item 综合题：每题 5--8 分钟；
  \item 判断题：先把 $x$ 换成 $-x$，再比较和原函数的关系。
\end{itemize}
\end{notebox}

\section{基础巩固}

\begin{exercise}
把函数
\[
y=x^2+1
\]
关于 x 轴对称，写出新方程。
\end{exercise}

\begin{exercise}
把函数
\[
y=2x-3
\]
关于 y 轴对称，写出新方程。
\end{exercise}

\begin{exercise}
把函数
\[
y=x^3
\]
关于原点对称，写出新方程。
\end{exercise}

\begin{exercise}
把函数
\[
y=|x-2|+1
\]
关于 y 轴对称，写出新方程。
\end{exercise}

\section{反向判断}

\begin{exercise}
已知函数
\[
y=-x^2+4,
\]
说出它是由
\[
y=x^2-4
\]
经过怎样的对称得到的。
\end{exercise}

\begin{exercise}
已知函数
\[
y=(x+3)^2,
\]
写出它关于 y 轴对称后的方程，并说明顶点怎样变化。
\end{exercise}

\begin{exercise}
已知函数
\[
y=x^4+2x^2+5,
\]
判断它的图像是否关于 y 轴对称，并说明理由。
\end{exercise}

\begin{exercise}
已知函数
\[
y=x^5-2x^3+x,
\]
判断它的图像是否关于原点对称，并说明理由。
\end{exercise}

\section{综合应用}

\begin{exercise}
把函数
\[
y=x^2-4x+1
\]
关于 y 轴对称，写出化简后的新方程。
\end{exercise}

\begin{exercise}
把函数
\[
y=\frac{1}{x-1}
\]
关于 y 轴对称，再关于 x 轴对称，写出最终方程。
\end{exercise}

\begin{exercise}
已知函数
\[
y=|x|
\]
经过某次对称后，顶点仍在原点，但图像方向改变了吗？请分别讨论关于 x 轴、关于 y 轴、关于原点对称时的结果。
\end{exercise}

\begin{exercise}
已知函数
\[
y=(x-2)^2+3,
\]
先关于 y 轴对称，再向下平移 5 个单位，写出最终方程。
\end{exercise}

\begin{exercise}
已知原图像上点 $(2,5)$ 在函数
\[
y=f(x)
\]
上。若图像关于 y 轴对称，写出对称后对应点的坐标，并写出新图像的方程。
\end{exercise}

\begin{exercise}
已知原图像上点 $(-3,4)$ 在函数
\[
y=f(x)
\]
上。若图像关于原点对称，写出对称后对应点的坐标，并写出新图像的方程。
\end{exercise}

\begin{exercise}
已知函数
\[
y=x^2+6x+8,
\]
说出它是由
\[
y=x^2-6x+8
\]
经过怎样的对称得到的。
\end{exercise}

\begin{exercise}
已知函数
\[
y=f(x)
\]
关于 x 轴对称后得到
\[
y=-f(x),
\]
再关于 y 轴对称，写出最终方程。
\end{exercise}

\newpage
\section{常见易错点总结}

\begin{goalbox}{做函数对称题时，最容易丢分的 5 个地方}
\begin{enumerate}
  \item \textbf{把关于 x 轴和关于 y 轴对称混掉。}  
  关于 x 轴是改外面，关于 y 轴是改里面。

  \item \textbf{关于 y 轴对称时，只改了一部分。}  
  必须把整个式子里的 $x$ 都换成 $-x$。

  \item \textbf{关于原点对称时漏了一个负号。}  
  正确形式是 $y=-f(-x)$。

  \item \textbf{不会用 $f(-x)=f(x)$ 或 $f(-x)=-f(x)$ 判断对称性。}  
  这是很多判断题最直接的方法。

  \item \textbf{不会用关键点检查。}  
  顶点、截距点、特殊点往往能快速发现结果是否合理。
\end{enumerate}
\end{goalbox}

\begin{notebox}{一个实用检查顺序}
\begin{itemize}
  \item 先判断是哪一种对称；
  \item 再决定是改“外面”还是改“里面”；
  \item 最后用特殊点或奇偶性检查结果。
\end{itemize}
\end{notebox}

\newpage
\appendix
\section{详细解答}

\subsection*{基础巩固}

\textbf{1}\begin{solution}
关于 x 轴对称，就是在原函数外面加负号，所以
\[
\boxed{y=-(x^2+1)=-x^2-1}.
\]
\end{solution}

\textbf{2}\begin{solution}
关于 y 轴对称，应把 $x$ 换成 $-x$：
\[
y=2(-x)-3=-2x-3.
\]
所以
\[
\boxed{y=-2x-3}.
\]
\end{solution}

\textbf{3}\begin{solution}
关于原点对称，应写成
\[
y=-f(-x).
\]
这里 $f(x)=x^3$，所以
\[
y=-(-x)^3=x^3.
\]
因此
\[
\boxed{y=x^3}.
\]
\end{solution}

\textbf{4}\begin{solution}
关于 y 轴对称，应把 $x$ 换成 $-x$：
\[
y=|(-x)-2|+1=|x+2|+1.
\]
所以
\[
\boxed{y=|x+2|+1}.
\]
\end{solution}

\subsection*{反向判断}

\textbf{5}\begin{solution}
原函数是
\[
y=x^2-4.
\]
关于 x 轴对称后应变成
\[
y=-(x^2-4)=-x^2+4.
\]
因此题目中的函数是由原函数\textbf{关于 x 轴对称}得到的。
\end{solution}

\textbf{6}\begin{solution}
关于 y 轴对称时，应把式子里的 $x$ 换成 $-x$：
\[
y=((-x)+3)^2=(x-3)^2.
\]
所以对称后的方程为
\[
\boxed{y=(x-3)^2}.
\]

原函数顶点是 $(-3,0)$，关于 y 轴对称后变成 $(3,0)$。
\end{solution}

\textbf{7}\begin{solution}
把 $x$ 换成 $-x$：
\[
f(-x)=(-x)^4+2(-x)^2+5=x^4+2x^2+5=f(x).
\]
因此图像\textbf{关于 y 轴对称}。
\end{solution}

\textbf{8}\begin{solution}
把 $x$ 换成 $-x$：
\[
f(-x)=(-x)^5-2(-x)^3+(-x)=-x^5+2x^3-x=-(x^5-2x^3+x)=-f(x).
\]
因此图像\textbf{关于原点对称}。
\end{solution}

\subsection*{综合应用}

\textbf{9}\begin{solution}
关于 y 轴对称，应把整个式子里的 $x$ 都换成 $-x$：
\[
y=(-x)^2-4(-x)+1=x^2+4x+1.
\]
所以
\[
\boxed{y=x^2+4x+1}.
\]
\end{solution}

\textbf{10}\begin{solution}
先关于 y 轴对称：
\[
y=\frac{1}{(-x)-1}=-\frac{1}{x+1}.
\]
再关于 x 轴对称：
\[
y=-\left(-\frac{1}{x+1}\right)=\frac{1}{x+1}.
\]
所以最终方程为
\[
\boxed{y=\frac{1}{x+1}}.
\]
\end{solution}

\textbf{11}\begin{solution}
原函数 $y=|x|$ 的图像是一个顶点在原点的 V 形。
\begin{itemize}
  \item 关于 x 轴对称后，变成倒着开的 V 形，方程是 $y=-|x|$；
  \item 关于 y 轴对称后，图像不变，方程仍是 $y=|x|$；
  \item 关于原点对称后，先左右翻，再上下翻，最终也是 $y=-|x|$。
\end{itemize}
\end{solution}

\textbf{12}\begin{solution}
先关于 y 轴对称：
\[
y=((-x)-2)^2+3=(x+2)^2+3.
\]
再向下平移 5 个单位：
\[
y=(x+2)^2+3-5=(x+2)^2-2.
\]
所以
\[
\boxed{y=(x+2)^2-2}.
\]
\end{solution}

\textbf{13}\begin{solution}
点 $(2,5)$ 关于 y 轴对称后，横坐标变号，纵坐标不变，所以对应点是
\[
\boxed{(-2,5)}.
\]

新图像的方程为
\[
\boxed{y=f(-x)}.
\]
\end{solution}

\textbf{14}\begin{solution}
点 $(-3,4)$ 关于原点对称后，横纵坐标都变号，所以对应点是
\[
\boxed{(3,-4)}.
\]

新图像的方程为
\[
\boxed{y=-f(-x)}.
\]
\end{solution}

\textbf{15}\begin{solution}
把
\[
y=x^2-6x+8
\]
关于 y 轴对称，应把式子里的 $x$ 换成 $-x$：
\[
y=(-x)^2-6(-x)+8=x^2+6x+8.
\]
所以题目中的函数是由原函数\textbf{关于 y 轴对称}得到的。
\end{solution}

\textbf{16}\begin{solution}
先关于 x 轴对称，得到
\[
y=-f(x).
\]
再关于 y 轴对称，应把其中的 $x$ 换成 $-x$，得到
\[
y=-f(-x).
\]
所以最终方程为
\[
\boxed{y=-f(-x)}.
\]
\end{solution}

\end{document}
